\apendice{Especificación de diseño}

\section{Introducción}

\section{Diseño de datos}

Para llevar a cabo la aplicación se han definido las siguientes entidades:
\begin{itemize}
\tightlist
	\item Usuario: esta entidad representa a la persona que usa la aplicación. Distingue entre usuario normal y administrador, al cual le permite realizar la gestión de documentos y usuarios. Esta entidad maneja un identificador (id) que es la clave primaria, un nombre, un email que debe ser único, una contraseña que se almacena \textit{hasheada}, el último inicio de sesión y el rol.
	\item Consulta: esta entidad guarda las preguntas que realizan los usuarios, asi como las respuestas que da el \textit{LLM} alimentado por el \textit{RAG}. Es la base para el historial de consultas donde se van a listar y borrar. Esta entidad maneja un identificador (id) que es la clave primaria, el identificador del usuario que realiza la consulta, la pregunta del usuario, la respuesta, los fragmentos utilizados, el tiempo de respuesta, y la fecha en la que se preguntó al modelo.
	\item Documento: esta entidad representa un pdf cargado en la aplicación. Maneja un identificador (id) como clave primaria, el nombre del fichero, el tamaño, la fecha de modificación, el número de \textit{chunks} en la base de datos vectorial, el \textit{hash} del documento y el estado del documento (cargado, procesado, indexado, fallido). 
	\item \textit{Chunk}: esta entidad representa cada trozo de texto en el que se dividen los documentos para poder generar los \textit{embeddings}. Maneja el texto que forma al propio \textit{chunk}, el \textit{embedding} y los metadatos del \textit{chunk} (nombre del archivo, título y posición dentro del documento).
	\item \textit{Embedding}: esta entidad se corresponde con la representación vectorial utilizada para la base de datos vectorial. Va a permitir recuperar \textit{chunks} por similitud. La entidad maneja un vector asociado al \textit{chunk}.
	\item \textit{ConsultaCunk}: esta entidad modela la relación entre las consultas y los \textit{chunks} permitiendo que una consulta tenga varios \textit{chunk} y que un \textit{chunk} pueda servir para varias consultas. La entidad maneja el identificador de la consulta, el identificador del \textit{chunk}, la similitud y el \textit{ranking}.
\end{itemize}

Las relaciones entre las entidades son:
\begin{itemize}
\tightlist
	\item Un usuario puede hacer cero o varias consultas.
	\item Cada consulta pertenece a un usuario.
	\item Un documento tiene cero (si esta cargado pero no indexado) o varios \textit{chunks}.
	\item Cada \textit{chunk} pertenece a un documento.
	\item Un \textit{chunk} puede tener cero o un \textit{embedding}.
	\item Cada \textit{embedding} pertenece a un \textit{chunk}.
	\item Una consulta usa uno o varios \textit{chunk}.
	\item Un \textit{chunk} puede ser usado por cero o varias consultas.
\end{itemize}


\imagenflotante{Diagramas/Entidad_Relacion}{Diagrama Entidad/Relación.}

\imagenflotante{Diagramas/Relacional}{Diagrama relacional.}


\section{Diseño arquitectónico}

\section{Diseño procedimental}

\section{Diseño de interfaces}



